\documentclass[11pt]{article}
\usepackage{amsmath,amssymb,amsthm}
\usepackage{mathtools}
\usepackage{tikz}
\usepackage{enumitem}
\usepackage[margin=1in]{geometry}

\newtheorem{theorem}{Theorem}
\newtheorem*{theorem*}{Theorem}
\newtheorem{lemma}{Lemma}
\newtheorem{definition}{Definition}
\newtheorem*{definition*}{Definition}
\newtheorem{proposition}{Proposition}
\newtheorem{assumption}{Assumption}
\newtheorem{remark}{Remark}

\title{Math 106 --Levy processes and rare events }
\author{Ethan Levien}
\date{}

\begin{document}
\maketitle


\section{Prelude: The generalized CLT}\label{sec:genclt}

Recall that Wiener processes can be understood as the scaling limit of a random walk where the steps have mean zero and finite variance. This result is a functional version of the CLT known as Donsker's theorem. 
L\'evy processes generalize Brownian motion by allowing step sizes with infinite variance (and, possibly, infinite mean). To understand this, we first discuss the generalized CLT for heavy-tailed increments. Power laws are representative of random variables with divergent moments and lead to stable limits.

Let $X_1,X_2,\dots$ be a sequence of i.i.d.\ random variables such that for some $\alpha \in (0,\infty)$
\begin{equation}\label{eq:tail_powerlaw}
\bar{F}(x) \coloneqq {\mathbb P}(X_1>x) \sim C x^{-\alpha},\quad \text{ as } x \to \infty.
\end{equation}
Therefore, the density of $X_1$ obeys 
\begin{equation}\label{eq:density_powerlaw}
f(x) \sim C\alpha\, x^{-\alpha-1},\quad x \to \infty.
\end{equation}
We assume the same tail behavior as $x\to -\infty$ with the \emph{same constant} $C$:
\begin{equation}\label{eq:two_sided_symmetric_tail}
{\mathbb P}(X_1>x)\sim C x^{-\alpha},
\qquad
{\mathbb P}(X_1<-x)\sim C x^{-\alpha},
\qquad x\to\infty.
\end{equation}
This symmetry assumption implies that the limiting stable law will be \emph{symmetric} (skewness parameter $\beta=0$). 

The mean does not exist if $\alpha\le 1$ and the variance does not exist if $\alpha\le 2$.

\begin{remark}[Mean in the symmetric case]
One might wonder whether the symmetry assumption \eqref{eq:two_sided_symmetric_tail} causes the divergences to cancel, making the mean well-defined even for $\alpha\le 1$. The answer is no, at least not in the standard sense. The mean is defined as ${\mathbb E}[X_1]=\int x\,f(x)\,dx$, which requires \emph{absolute} integrability: ${\mathbb E}[|X_1|]=\int_{0}^{\infty}{\mathbb P}(|X_1|>x)\,dx\asymp \int_1^\infty x^{-\alpha}\,dx$. This integral diverges when $\alpha\le 1$, regardless of symmetry.

What symmetry \emph{does} give is that the \emph{principal-value} integral $\lim_{R\to\infty}\int_{-R}^{R}x\,f(x)\,dx=0$ is well-defined and equal to zero, since the integrand is odd. This is why no centering ($A_n=0$) is needed in the generalized CLT even for $\alpha<1$: the characteristic function is real and even by symmetry, so the limiting law is automatically symmetric around zero. But this is a property of symmetry, not of the mean existing.
\end{remark}

We study the partial sums
\begin{equation}\label{eq:Sn_def}
S_n = \sum_{i=1}^n X_i.
\end{equation}

\subsection{Goal and characteristic functions}\label{subsec:goal_cf}

Motivated by the classical CLT, we ask: can we find a normalization $B_n\to\infty$ such that
\begin{equation}\label{eq:goal_Sn}
\frac{S_n}{B_n}\;\Rightarrow\; Z
\quad \text{as } n\to\infty,
\end{equation}
for some non-degenerate limit random variable $Z$?

It is easier to work with characteristic functions. Define
\begin{equation}\label{eq:phi_def}
\phi(t)\coloneqq {\mathbb E}\!\left[e^{-itX_1}\right],
\end{equation}
and
\begin{equation}\label{eq:phi_n_def}
\phi_n(u) \coloneqq {\mathbb E}\!\left[e^{-iuS_n/B_n}\right]
= \left({\mathbb E}\!\left[e^{-iuX_1/B_n}\right]\right)^n
= \left(\phi(u/B_n)\right)^n.
\end{equation}
Thus the problem reduces to understanding the small-$t$ asymptotics of $\phi(t)$ as $t\to 0$.

\subsection{A Tauberian theorem in the symmetric power-law case}\label{subsec:tauberian_symmetric}

The link between power-law tails and small-$t$ characteristic functions is a Tauberian theorem. In our symmetric setting \eqref{eq:two_sided_symmetric_tail}, it takes a particularly clean form.

\begin{theorem}[Tauberian theorem (symmetric power-law tails $\leftrightarrow$ fractional characteristic exponent)]\label{thm:tauberian_symmetric}
Fix $\alpha\in(0,2)$ with $\alpha\neq 1$. Assume the two-sided symmetric tail condition \eqref{eq:two_sided_symmetric_tail}. Then as $t\to 0$,
\begin{equation}\label{eq:tauberian_forward}
1-\phi(t)
\;=\;
2C\,\Gamma(1-\alpha)\,\cos\!\Big(\frac{\pi\alpha}{2}\Big)\,|t|^\alpha
\;+\; o(|t|^\alpha).
\end{equation}
Conversely, if for some $K>0$ one has
\begin{equation}\label{eq:tauberian_converse_assump}
1-\phi(t) = K|t|^\alpha + o(|t|^\alpha)
\quad \text{as } t\to 0,
\end{equation}
then $X_1$ has two-sided regularly varying tails of index $\alpha$ and, in the symmetric case,
\begin{equation}\label{eq:tauberian_converse_tail}
{\mathbb P}(X_1>x)\sim C x^{-\alpha},
\qquad
{\mathbb P}(X_1<-x)\sim C x^{-\alpha},
\qquad x\to\infty,
\end{equation}
with
\begin{equation}\label{eq:C_from_K}
C \;=\; \frac{K}{2\,\Gamma(1-\alpha)\,\cos(\pi\alpha/2)}.
\end{equation}
\end{theorem}

\subsection{Sketch proof of Theorem \ref{thm:tauberian_symmetric}}\label{subsec:tauberian_proof_sketch}

We sketch the forward direction \eqref{eq:tauberian_forward} under \eqref{eq:two_sided_symmetric_tail}; the converse is a reverse Tauberian argument using the same integral identities.

Because the tails are symmetric, $\phi(t)$ is real and even:
\begin{equation}\label{eq:phi_real_even}
\phi(t)= {\mathbb E}[\cos(tX_1)],
\qquad
1-\phi(t) = {\mathbb E}[1-\cos(tX_1)].
\end{equation}
By symmetry $\phi(t)={\mathbb E}[\cos(tX_1)]$, so $1-\phi(t)={\mathbb E}[1-\cos(tX_1)]$. Applying integration by parts: for any nonnegative, smooth $g$ with $g(0)=0$,
\begin{equation}\label{eq:ibp_template}
{\mathbb E}[g(|X_1|)]
= \int_0^\infty g'(x)\,{\mathbb P}(|X_1|>x)\,dx.
\end{equation}
Apply \eqref{eq:ibp_template} with $g(x)=1-\cos(tx)$ (note $g(0)=0$ and $g'(x)=t\sin(tx)$):
\begin{equation}\label{eq:ibp_cos}
1-\phi(t)
= \int_0^\infty t\sin(tx)\,{\mathbb P}(|X_1|>x)\,dx.
\end{equation}
Under \eqref{eq:two_sided_symmetric_tail}, we have ${\mathbb P}(|X_1|>x)\sim 2C x^{-\alpha}$, so heuristically
\begin{equation}\label{eq:asympt_substitution}
1-\phi(t)
\sim \int_0^\infty t\sin(tx)\, (2C x^{-\alpha})\,dx.
\end{equation}
Change variables $y=tx$:
\begin{equation}\label{eq:change_var}
\int_0^\infty t\sin(tx)\,x^{-\alpha}\,dx
= |t|^\alpha \int_0^\infty \sin y \, y^{-\alpha}\,dy,
\end{equation}
where the power of $|t|$ comes from $x^{-\alpha}=(y/|t|)^{-\alpha}$ and $dx=dy/|t|$.
Thus we expect
\begin{equation}\label{eq:1_minus_phi_constant}
1-\phi(t)\sim 2C |t|^\alpha \int_0^\infty \sin y\, y^{-\alpha}\,dy.
\end{equation}
The remaining integral is classical and can be evaluated via analytic continuation of the Gamma function:
\begin{equation}\label{eq:sin_integral}
\int_0^\infty \sin y\, y^{-\alpha}\,dy
= \Gamma(1-\alpha)\cos\!\Big(\frac{\pi\alpha}{2}\Big),
\qquad \alpha\in(0,2),\ \alpha\neq 1,
\end{equation}
which yields exactly \eqref{eq:tauberian_forward}.

To justify replacing ${\mathbb P}(|X_1|>x)$ by its asymptotic $2C x^{-\alpha}$ inside \eqref{eq:ibp_cos}, one splits the integral at $x=1/|t|$:
\begin{equation}\label{eq:split_scale}
\int_0^\infty t\sin(tx)\,{\mathbb P}(|X_1|>x)\,dx
=
\int_0^{1/|t|} (\cdots)\,dx \;+\; \int_{1/|t|}^{\infty}(\cdots)\,dx,
\end{equation}
uses that $\sin(tx)\approx tx$ on $[0,1/|t|]$ together with regular variation to control the first piece, and uses boundedness of $\sin$ plus the power-law tail to control the second piece. The $o(|t|^\alpha)$ remainder comes from the fact that the ratio ${\mathbb P}(|X_1|>x)/(2C x^{-\alpha})\to 1$ slowly but uniformly on the relevant $x$-scales after the change of variables $y=tx$.

\subsection{Heuristic derivation of the (symmetric) stable characteristic function}\label{subsec:stable_cf_heuristic}

Assume $\alpha\in(0,2)$, $\alpha\neq 1$, and the symmetric tail condition \eqref{eq:two_sided_symmetric_tail}. By Theorem \ref{thm:tauberian_symmetric},
\begin{equation}\label{eq:small_t_phi_symmetric}
\phi(t)
=
1 - K|t|^\alpha + o(|t|^\alpha),
\qquad
K \coloneqq 2C\,\Gamma(1-\alpha)\cos\!\Big(\frac{\pi\alpha}{2}\Big).
\end{equation}
Plugging $t=u/B_n$ into \eqref{eq:phi_n_def} gives
\begin{equation}\label{eq:phi_expand}
\phi(u/B_n)
=
1 - K|u|^\alpha B_n^{-\alpha} + o(B_n^{-\alpha}).
\end{equation}
Therefore
\begin{align}
\log \phi_n(u)
&= n\log(\phi(u/B_n))
\label{eq:log_phi_n_1}\\
&= n\log\Big(1 - K|u|^\alpha B_n^{-\alpha} + o(B_n^{-\alpha})\Big)
\label{eq:log_phi_n_2}\\
&= -nK|u|^\alpha B_n^{-\alpha} + o\!\left(nB_n^{-\alpha}\right),
\label{eq:log_phi_n_3}
\end{align}
using $\log(1-z)=-z+o(z)$. Choosing the normalization
\begin{equation}\label{eq:Bn_choice}
B_n = n^{1/\alpha},
\end{equation}
we obtain $nB_n^{-\alpha}=1$ and hence
\begin{equation}\label{eq:phi_n_limit}
\phi_n(u)\;\longrightarrow\;\exp\!\left(-K|u|^\alpha\right).
\end{equation}
The limit characteristic function \eqref{eq:phi_n_limit} defines the \emph{symmetric $\alpha$-stable} law (Gaussian when $\alpha=2$). This is the heavy-tailed analogue of the classical CLT: after normalization by $n^{1/\alpha}$, the sum converges to a universal family indexed by $\alpha$ (and, in the nonsymmetric case, by an additional skewness parameter).

\subsection{The limit law has a power-law tail of the same index}\label{subsec:stable_tail}

The same Tauberian theorem also explains why the stable limit has a power-law tail of index $\alpha$. Let $Z$ have characteristic function
\begin{equation}\label{eq:stable_cf_symmetric}
{\mathbb E}[e^{-iuZ}] = \exp(-K|u|^\alpha),
\end{equation}
so that for small $u$,
\begin{equation}\label{eq:stable_small_u}
1-\phi_Z(u) = K|u|^\alpha + o(|u|^\alpha).
\end{equation}
Applying the \emph{converse} direction in Theorem \ref{thm:tauberian_symmetric} yields
\begin{equation}\label{eq:stable_tail_conclusion}
{\mathbb P}(Z>x)\sim \widetilde C x^{-\alpha},
\qquad
{\mathbb P}(Z<-x)\sim \widetilde C x^{-\alpha},
\qquad x\to\infty,
\end{equation}
with $\widetilde C$ related to $K$ by \eqref{eq:C_from_K}. In particular, the tail exponent $\alpha$ is preserved by the scaling limit.

\section{Scaling and extremes: why $B_n\sim n^{1/\alpha}$}\label{sec:scaling_extremes}

The normalization $B_n=n^{1/\alpha}$ can also be motivated directly from the order of the maximum. Let
\begin{equation}\label{eq:max_def}
M_n \coloneqq \max_{1\le i\le n} |X_i|.
\end{equation}
Under \eqref{eq:two_sided_symmetric_tail},
\begin{equation}\label{eq:abs_tail}
{\mathbb P}(|X_1|>x) \sim 2C x^{-\alpha}.
\end{equation}
For any fixed $x>0$,
\begin{align}
{\mathbb P}\!\left(M_n \le n^{1/\alpha}x\right)
&=
{\mathbb P}\!\left(|X_1|\le n^{1/\alpha}x\right)^n
=
\left(1-{\mathbb P}\!\left(|X_1|>n^{1/\alpha}x\right)\right)^n
\label{eq:max_cdf_1}\\
&\approx
\left(1- 2C (n^{1/\alpha}x)^{-\alpha}\right)^n
=
\left(1-\frac{2C}{n}x^{-\alpha}\right)^n
\label{eq:max_cdf_2}\\
&\longrightarrow \exp(-2C x^{-\alpha})
\qquad (n\to\infty).
\label{eq:max_limit}
\end{align}
Thus $M_n$ is typically of order $n^{1/\alpha}$. Since the dominant contribution to $S_n$ comes from the largest terms, which are of order $n^{1/\alpha}$, it is natural to normalize by $B_n = n^{1/\alpha}$.



\section{Generalized CLT (Fisher--Gnedenko / Gnedenko--Kolmogorov)}\label{sec:gclt_statement}

An $\alpha$-stable random variable $Z$ ($0<\alpha\le 2$) is characterized by stability under summation: for independent copies $Z_1,Z_2$ with the same law as $Z$,
\begin{equation}\label{eq:stability_property}
Z_1+Z_2 \stackrel{d}{=} 2^{1/\alpha} Z
\end{equation}
(in the symmetric, zero-drift case). This follows directly from the characteristic function: ${\mathbb E}[e^{iu(Z_1+Z_2)}] = e^{-2\sigma^\alpha|u|^\alpha} = {\mathbb E}[e^{iu\cdot 2^{1/\alpha}Z}]$.

In the symmetric case relevant to our standing assumption \eqref{eq:two_sided_symmetric_tail}, the characteristic function takes the form
\begin{equation}\label{eq:symmetric_stable_cf}
{\mathbb E}\!\left[e^{iuZ}\right] = \exp\!\left(-\sigma^\alpha |u|^\alpha\right),
\qquad 0<\alpha\le 2,
\end{equation}
for some scale parameter $\sigma>0$. When $\alpha=2$ this is Gaussian. When $\alpha<2$, $Z$ has power-law tails of the same index:
\begin{equation}\label{eq:stable_tail_again}
{\mathbb P}(Z>z)\sim \widetilde C z^{-\alpha},
\qquad z\to\infty,
\end{equation}
(and symmetrically for $Z<-z$).
More generally (dropping symmetry), one introduces an additional skewness parameter $\beta\in[-1,1]$ in the exponent; we omit this here since our tails are symmetric.


The full theorem classifies all possible non-degenerate scaling limits of i.i.d.\ sums. In our notes we only need the heavy-tail stable cases and the classical CLT as a comparison.

\begin{theorem}[Generalized CLT / Fisher--Gnedenko theorem]\label{thm:gclt}
Let $X_1,X_2,\dots$ be i.i.d.\ and set $S_n=\sum_{i=1}^n X_i$. Suppose there exist sequences $A_n\in\mathbb R$ and $B_n>0$ such that
\begin{equation}\label{eq:gclt_conv}
\frac{S_n-A_n}{B_n}\Rightarrow \eta
\quad \text{as } n\to\infty,
\end{equation}
where $\eta$ is non-degenerate. Then $\eta$ is either Gaussian (the classical CLT case) or an $\alpha$-stable law for some $\alpha\in(0,2]$.

Moreover, if $X_1$ has power-law tails of index $\alpha\in(0,2)$ then one may take $B_n$ of order $n^{1/\alpha}$, and $\eta$ is $\alpha$-stable (symmetric under \eqref{eq:two_sided_symmetric_tail}). If $X_1$ has finite variance, then one may take $B_n$ of order $\sqrt{n}$ and $\eta$ is Gaussian.
\end{theorem}

A convenient summary (in our notation, with symmetry so $A_n=0$ in the $\alpha<2$ rows) is:

\begin{table}[h]
\centering
\begin{tabular}{lllll}
\hline
Tail index $\alpha$ & ${\mathbb E}[X_1]$ & ${\rm var}(X_1)$ & Typical $B_n$ & Limit law $\eta$ \\
\hline
$2<\alpha$ (finite variance) & finite & finite & $B_n=O(\sqrt{n})$ & Standard Normal \\
$1<\alpha<2$ & finite & $\infty$ & $B_n=O(n^{1/\alpha})$ & $\alpha$-stable \\
$0<\alpha<1$ & $\infty$ & $\infty$ & $B_n=O(n^{1/\alpha})$ & $\alpha$-stable \\
\hline
\end{tabular}
\caption{Generalized CLT scaling regimes. Under our symmetric tail assumption, the stable limits are symmetric.}
\label{tab:gclt}
\end{table}


\section{From random walks to L\'evy processes via heavy-tailed invariance principles}\label{sec:rw_to_levy}

\subsection{L\'evy processes and symmetric $\alpha$-stable processes}\label{subsec:levy_def_stable}

\begin{definition}[L\'evy process]\label{def:levy_process}
A stochastic process $\{L_t\}_{t\ge 0}$ is a \emph{L\'evy process} if:
\begin{enumerate}[label=(\roman*)]
\item $L_0=0$ almost surely;
\item $L$ has \emph{stationary increments}: for $0\le s<t$, $L_t-L_s$ has the same law as $L_{t-s}$;
\item $L$ has \emph{independent increments}: for $0\le t_0<t_1<\cdots<t_m$, the increments $L_{t_1}-L_{t_0},\dots,L_{t_m}-L_{t_{m-1}}$ are independent;
\item $L$ is \emph{stochastically continuous}: for every $t\ge 0$, $L_{t+h}\to L_t$ in probability as $h\to 0$.
\end{enumerate}
\end{definition}

A L\'evy process admits a modification with c\`adl\`ag sample paths (right-continuous with left limits). In general, such paths need not be continuous and typically have jumps.

\begin{definition}[Symmetric $\alpha$-stable L\'evy process]\label{def:saslp}
Fix $\alpha\in(0,2]$ and a scale parameter $\sigma>0$. A L\'evy process $\{Z_t\}_{t\ge 0}$ is called a \emph{symmetric $\alpha$-stable process} if for each $t\ge 0$,
\begin{equation}\label{eq:sas_char}
{\mathbb E}\!\left[e^{iuZ_t}\right]
= \exp\!\left(-t\,\sigma^\alpha |u|^\alpha\right),
\qquad u\in\mathbb R.
\end{equation}
Equivalently, $Z_t$ is $\alpha$-stable for each $t$ and the increments are stable under addition:
\begin{equation}\label{eq:stable_increments}
Z_t-Z_s \;\stackrel{d}{=}\; Z_{t-s},
\qquad 0\le s<t,
\end{equation}
with the stability property that for independent copies $Z^{(1)}_{t},Z^{(2)}_{t}$,
\begin{equation}\label{eq:stability_addition}
Z^{(1)}_{t}+Z^{(2)}_{t} \;\stackrel{d}{=}\; 2^{1/\alpha} Z_{t}.
\end{equation}
\end{definition}

When $\alpha=2$, \eqref{eq:sas_char} is Gaussian and $Z$ is Brownian motion (up to variance normalization). When $\alpha\in(0,2)$, $Z$ has jumps and heavy-tailed increments.

\subsection{Why the limiting sample paths are no longer continuous (heuristic)}\label{subsec:paths_not_continuous}

Consider the random walk partial sums $S_n=\sum_{k=1}^n X_k$ with i.i.d.\ increments satisfying the symmetric tail assumption
\begin{equation}\label{eq:symm_tail_again}
{\mathbb P}(|X_1|>x)\sim 2C x^{-\alpha},
\qquad x\to\infty,
\end{equation}
for some $\alpha\in(0,2)$. Define the linearly time-rescaled process
\begin{equation}\label{eq:scaled_rw_process}
X^{(n)}(t)\;\coloneqq\;\frac{1}{n^{1/\alpha}}\sum_{k=1}^{\lfloor nt\rfloor} X_k,
\qquad t\in[0,T].
\end{equation}
The generalized CLT suggests that $\{X^{(n)}(t)\}$ converges (in $D([0,T])$) to a symmetric $\alpha$-stable L\'evy process $\{Z_t\}$.

The key qualitative point is that, under the $n^{1/\alpha}$ scaling, \emph{macroscopic jumps persist in the limit}. Fix $\varepsilon>0$ and consider the number of increments among the first $\lfloor nt\rfloor$ steps whose \emph{rescaled} size exceeds $\varepsilon$:
\begin{equation}\label{eq:big_jumps_count}
N^{(n)}_{\varepsilon}(t)\;\coloneqq\;\sum_{k=1}^{\lfloor nt\rfloor}\mathbf 1\{|X_k|>\varepsilon n^{1/\alpha}\}.
\end{equation}
Then
\begin{equation}\label{eq:mean_big_jumps}
{\mathbb E}[N^{(n)}_{\varepsilon}(t)]
=
\lfloor nt\rfloor\,{\mathbb P}(|X_1|>\varepsilon n^{1/\alpha})
\sim
nt \cdot 2C (\varepsilon n^{1/\alpha})^{-\alpha}
=
2Ct\,\varepsilon^{-\alpha}.
\end{equation}
Thus the expected number of $\varepsilon$-jumps over a fixed time horizon is of order $1$, not of order $o(1)$ as in the finite-variance (Brownian) scaling. In fact, a standard rare-event heuristic says
\begin{equation}\label{eq:poisson_heuristic}
N^{(n)}_{\varepsilon}(t) \;\approx\; \mathrm{Poisson}(2Ct\,\varepsilon^{-\alpha}),
\end{equation}
so
\begin{equation}\label{eq:jump_probability}
{\mathbb P}\!\left(\sup_{0\le s\le t} |X^{(n)}(s)-X^{(n)}(s-)|>\varepsilon\right)
=
{\mathbb P}\!\left(N^{(n)}_{\varepsilon}(t)\ge 1\right)
\;\to\;
1-\exp(-2Ct\,\varepsilon^{-\alpha}).
\end{equation}
The right-hand side is strictly positive for every $t>0$. Therefore, in the limit one cannot have continuous sample paths: with positive probability there is a jump of size at least $\varepsilon$ on $[0,t]$, for every $\varepsilon>0$.

This contrasts with the Brownian case: under $\sqrt{n}$ scaling, the probability of an $O(1)$ rescaled jump vanishes because ${\mathbb P}(|X_1|>\varepsilon\sqrt{n})$ is too small when ${\mathbb E}[X_1^2]<\infty$.

\subsection{Generator of a symmetric $\alpha$-stable process}\label{subsec:stable_generator}

\paragraph{The fractional Laplacian.}
Before identifying the generator, we define the key operator that will appear.

\begin{definition}[Fractional Laplacian]\label{def:frac_lap}
For $\beta>0$, the \emph{fractional Laplacian} $(-\Delta)^{\beta}$ is the operator defined on the Schwartz class (or $L^2$) by its Fourier multiplier:
\begin{equation}\label{eq:frac_lap_def}
\widehat{(-\Delta)^{\beta} f}(\xi) \;=\; |\xi|^{2\beta}\,\hat f(\xi).
\end{equation}
The notation $(-\Delta)^{\beta}$ reflects that $-\Delta$ has Fourier symbol $|\xi|^2$, so raising it to the power $\beta$ replaces $|\xi|^2$ by $|\xi|^{2\beta}$. The case $\beta=1/2$ gives $(-\Delta)^{1/2}$ with symbol $|\xi|$; more generally $(-\Delta)^{\alpha/2}$ has symbol $|\xi|^\alpha$.
\end{definition}

\paragraph{Deriving the generator from the semigroup.}

A symmetric $\alpha$-stable process is a (translation-invariant) Markov process, so it possesses a generator. Recall that the generator $\mathcal{A}$ of a semigroup $\{P_t\}_{t\ge 0}$ is defined as the operator
\begin{equation}\label{eq:generator_def}
\mathcal{A}f \;\coloneqq\; \lim_{t\downarrow 0}\frac{P_t f - f}{t},
\end{equation}
for functions $f$ in the domain of $\mathcal{A}$, where $P_t f(x) = {\mathbb E}[f(x + Z_t)]$ is the \emph{Markov semigroup}:
\begin{equation}\label{eq:stable_semigroup}
(P_t f)(x)\;\coloneqq\;{\mathbb E}\!\left[f(x+Z_t)\right].
\end{equation}
To identify $\mathcal{A}$, we work in the Fourier domain. Taking the Fourier transform of $P_t f$ and using the characteristic function \eqref{eq:sas_char}:
\begin{align}
\widehat{P_t f}(\xi)
&= \int_{\mathbb R} e^{-i\xi x}\,{\mathbb E}[f(x+Z_t)]\,dx \notag\\
&= {\mathbb E}\!\left[\int_{\mathbb R} e^{-i\xi x} f(x+Z_t)\,dx\right] \notag\\
&= {\mathbb E}\!\left[e^{i\xi Z_t}\right]\hat f(\xi)
\;=\;
e^{-t\sigma^\alpha|\xi|^\alpha}\,\hat f(\xi).
\label{eq:fourier_symbol}
\end{align}
(The change of variables $x \mapsto x - Z_t$ yields the factor ${\mathbb E}[e^{i\xi Z_t}]$, which is the characteristic function of $Z_t$ at $\xi$; by \eqref{eq:sas_char} this equals $e^{-t\sigma^\alpha|\xi|^\alpha}$.) Now apply the definition \eqref{eq:generator_def} in the Fourier domain:
\begin{equation}
\widehat{\mathcal{A}f}(\xi)
= \lim_{t\downarrow 0}\frac{\widehat{P_t f}(\xi)-\hat f(\xi)}{t}
= \lim_{t\downarrow 0}\frac{e^{-t\sigma^\alpha|\xi|^\alpha}-1}{t}\,\hat f(\xi)
= -\sigma^\alpha|\xi|^\alpha\,\hat f(\xi).
\label{eq:generator_symbol}
\end{equation}
Comparing with Definition \ref{def:frac_lap} (with $2\beta=\alpha$, i.e.\ $\beta=\alpha/2$), we recognize this as the fractional Laplacian:
\begin{equation}\label{eq:frac_lap}
\mathcal{A}f(x) = -\sigma^\alpha\,(-\Delta)^{\alpha/2} f(x).
\end{equation}

Equivalently, $\mathcal{A}$ can be written as an integro-differential operator via the L\'evy measure
\begin{equation}\label{eq:stable_levy_measure}
\nu(dy) = \frac{\kappa_\alpha\,\sigma^\alpha}{|y|^{1+\alpha}}\,dy,
\qquad y\in\mathbb R\setminus\{0\},
\end{equation}
for an explicit constant $\kappa_\alpha>0$ (its exact value depends on the normalization convention for $(-\Delta)^{\alpha/2}$). For sufficiently smooth $f$ (e.g.\ $f\in C_c^2(\mathbb R)$),
\begin{equation}\label{eq:generator_levy_form}
\mathcal{A}f(x)
=
\int_{\mathbb R\setminus\{0\}}
\Big(
f(x+y)-f(x)-y f'(x)\mathbf 1_{\{|y|<1\}}
\Big)\,\nu(dy).
\end{equation}
The subtraction term is needed precisely when $\int_{|y|<1}|y|\,\nu(dy)=\infty$, i.e.\ when $\alpha\ge 1$:
\begin{equation}\label{eq:compensation_condition}
\int_{|y|<1}|y|\,\nu(dy)
\asymp \int_0^1 r\cdot r^{-1-\alpha}\,dr
=
\int_0^1 r^{-\alpha}\,dr
=
\begin{cases}
\frac{1}{1-\alpha}, & \alpha<1,\\
\infty, & \alpha\ge 1.
\end{cases}
\end{equation}
Hence for $\alpha\in(0,1)$, one may omit the compensation term and simply write
\begin{equation}\label{eq:generator_alpha_lt_1}
\mathcal{A}f(x)
=
\int_{\mathbb R\setminus\{0\}}
\big(f(x+y)-f(x)\big)\,\nu(dy),
\qquad \alpha\in(0,1).
\end{equation}

\subsection{From compound Poisson to stable processes, and beyond}\label{subsec:compound_poisson_to_levy}

The jump structure above can be motivated by approximating stable processes using compound Poisson processes, and this also motivates the general L\'evy--Khintchine representation.

\paragraph{Step 1: compound Poisson approximation by truncating small jumps.}
Fix $\varepsilon>0$ and consider the truncated L\'evy measure
\begin{equation}\label{eq:truncated_levy_measure}
\nu_\varepsilon(dy)\;\coloneqq\;\mathbf 1_{\{|y|>\varepsilon\}}\,\nu(dy).
\end{equation}
Since $\nu_\varepsilon(\mathbb R)<\infty$, there is a compound Poisson process $Z^{(\varepsilon)}$ with jump intensity
\begin{equation}\label{eq:lambda_eps}
\lambda_\varepsilon \;\coloneqq\;\nu_\varepsilon(\mathbb R)
=
\int_{|y|>\varepsilon}\nu(dy)
\asymp \varepsilon^{-\alpha},
\end{equation}
and jump distribution proportional to $\nu_\varepsilon$, such that
\begin{equation}\label{eq:compound_poisson_form}
Z^{(\varepsilon)}_t = \sum_{j=1}^{N_t^{(\varepsilon)}} Y_j^{(\varepsilon)},
\end{equation}
where $N_t^{(\varepsilon)}\sim \mathrm{Poisson}(\lambda_\varepsilon t)$ and the jumps $Y_j^{(\varepsilon)}$ are i.i.d.\ with law $\nu_\varepsilon(dy)/\lambda_\varepsilon$.

This process has only finitely many jumps on compact time intervals, hence piecewise-constant paths. As $\varepsilon\downarrow 0$, the intensity $\lambda_\varepsilon\to\infty$ and one sees more and more small jumps; the limiting object has infinitely many small jumps on every interval (infinite activity), yet still has c\`adl\`ag paths.

\paragraph{Step 2: taking $\varepsilon\downarrow 0$ recovers the stable process.}
Heuristically, the family $Z^{(\varepsilon)}$ converges (after adding the compensation term when $\alpha\ge 1$) to the full symmetric $\alpha$-stable process $Z$. This is the simplest instance of the general L\'evy--It\^o decomposition: a L\'evy process can be built by superposing jumps across all scales according to a L\'evy measure $\nu$.

\paragraph{Step 3: motivation for general L\'evy processes and the L\'evy--Khintchine formula.}
The stable case corresponds to the special choice $\nu(dy)\propto |y|^{-1-\alpha}dy$. If instead we allow an arbitrary jump intensity across sizes, plus a Gaussian component and a deterministic drift, we obtain the general class of L\'evy processes.

The characteristic exponent of a general L\'evy process has the L\'evy--Khintchine form:
\begin{equation}\label{eq:levy_khintchine}
{\mathbb E}\!\left[e^{iuL_t}\right]
=
\exp\!\left(
t\,\Psi(u)
\right),
\qquad
\Psi(u)
=
i b u
-\frac{1}{2}\sigma^2 u^2
+
\int_{\mathbb R\setminus\{0\}}
\left(e^{iuy}-1-iuy\,\mathbf 1_{\{|y|<1\}}\right)\nu(dy),
\end{equation}
where $b\in\mathbb R$ is a drift, $\sigma\ge 0$ is a Brownian volatility, and $\nu$ is a L\'evy measure satisfying
\begin{equation}\label{eq:levy_measure_condition}
\int_{\mathbb R\setminus\{0\}} (1\wedge y^2)\,\nu(dy) < \infty.
\end{equation}
In this representation, the integral term encodes the aggregate effect of jumps of all sizes, with the same compensation mechanism as in \eqref{eq:generator_levy_form}. The symmetric $\alpha$-stable process is recovered by setting $b=0$, $\sigma=0$ (no Brownian component), and $\nu(dy)=c_\alpha |y|^{-1-\alpha}dy$ for an appropriate constant $c_\alpha>0$.


\end{document}